This paper developed a resilient estimation architecture that combines a Kalman filter, chi-square innovation monitoring, physics-guided attack generation, and channel-wise localization. The attack generator was formulated using finite-horizon maps that retain the dynamics of both the estimator and the controlled plant. This construction gives a detector-calibrated measure of stealthiness through the noncentral chi-square distribution, rather than through an unnormalized residual threshold.

The localization output was converted into an estimator input through calibrated safety probabilities. Under explicit calibration and conditional-independence assumptions, the selected subset is attack-free with a prescribed probability. The analysis further showed that attack sparsity limits the existence of a clean observable subset, while the size of the selected subset determines the conditional Kalman-filter covariance through the Riccati equation. The water-distribution case study demonstrated the empirical sensing redundancy and the localization improvement obtained by augmenting the training set with physics-guided attacks. Closed-loop numerical validation of the revised KF/$\chi^2$ architecture is outside the scope of the available data and is not inferred from the earlier residual-based results.

Future work will consider dependent channel-calibration models, time-varying sensor subsets, and unconditional error guarantees obtained by combining probabilistic selection with a robust update on the selection-failure event.
